\section{What the Body Knows}

The first of religion's own acts is invisible, and Aquinas gives it a name that has narrowed sadly in modern mouths: \emph{devotion}. The word now suggests a temperature---warmth of feeling, a candle's worth of sentiment. Its root is \latin{devovere}, to vow away, and Aquinas's definition keeps the root's edge: devotion is ``the will to give oneself readily to things concerning the service of God.''\endnote{\emph{Summa theologiae} II-II, q.~82, a.~1, corpus (\latin{devotio nihil aliud esse videtur quam voluntas quaedam prompte tradendi se ad ea quae pertinent ad Dei famulatum}); checked 2026-07-24, English at New Advent, Latin at Corpus Thomisticum. The article's opening ties the noun to \latin{devovere}: the devout are those who ``in a way, devote themselves to God, so as to subject themselves wholly to Him.''} A \emph{will}, not a feeling; and its differentia is not intensity but \emph{promptness}---\latin{prompte}, readily, without the interior dragging of feet. Devotion is what consent (Part I's word) becomes when it stops being a single act and becomes a disposition: the will's standing readiness to be handed over. Everything else religion does is this readiness taking shape in matter.

Its first shape is speech. And here the analysis meets the oldest objection to petitionary prayer, which Aquinas states as well as any contemporary: God knows the need before the asking; God's disposition does not change; therefore prayer informs no one and alters nothing; therefore it is empty. The objection assumes---note the family resemblance to recoil's mistake in Part I---that prayer, to be real, would have to act \emph{on God}: push the divine will as one billiard ball pushes another, and since it cannot, it must be theater. Aquinas's answer relocates the whole transaction inside the created order: ``we pray not that we may change the Divine disposition, but that we may impetrate that which God has disposed to be fulfilled by our prayers.''\endnote{\emph{Summa theologiae} II-II, q.~83, a.~2, corpus (\latin{non enim propter hoc oramus ut divinam dispositionem immutemus, sed ut id impetremus quod Deus disposuit per orationes sanctorum esse implendum}), quoting Gregory the Great's \emph{Dialogues}: ``that by asking, men may deserve to receive what Almighty God from eternity has disposed to give.'' Checked 2026-07-24 as above. The underlying scheme is the primary/secondary causality summarized in Part I from ST I, q.~105, a.~5: divine providence disposes not only which effects occur but through which causes they occur.} Providence disposes not only what will be given but \emph{through what} it will be given, and among the throughs are asking wills. Prayer is a secondary cause among secondary causes: as the eye sees because light is poured into it, the prayer obtains because giving is disposed through it. Nothing pushes upward; everything is received downward, including the asking itself; and the asking is nevertheless really causal, in exactly the way Part I said creaturely action is really action. A person praying is dependence doing what dependence is for: not informing God but \emph{practicing the truth}---taking up, on purpose, the position of a receiver, which was the true position all along.

Then the body itself enters, and the tradition is blunt about why it must. Man is not an angel inconvenienced by flesh; he is ``composed of a twofold nature, intellectual and sensible,'' and therefore ``we offer God a twofold adoration; namely, a spiritual adoration, consisting in the internal devotion of the mind; and a bodily adoration, which consists in an exterior humbling of the body.'' The direction of service runs inward: ``we exhibit signs of humility in our bodies in order to incite our affections to submit to God, since it is connatural to us to proceed from the sensible to the intelligible.''\endnote{\emph{Summa theologiae} II-II, q.~84, a.~2, corpus, citing John Damascene, \emph{De fide orthodoxa} IV.12; checked 2026-07-24 as above. Part I cited q.~84, a.~1 (adoration owed to God on account of the divine excellence) in its closing paragraph; the present section is that citation unfolded. On the priority of the interior act see also II-II, q.~81, a.~7: exterior worship exists for interior worship, not conversely.} The knee does not bend because God enjoys watching knees; it bends because the one kneeling learns from his own posture, the way a child learns courage by standing up straight. Aquinas's gloss on the extreme case compresses the whole of Part I into six words: when we genuflect, he says, we signify our weakness in comparison with God; but ``when we prostrate ourselves we profess that we are nothing of ourselves.''\endnote{\emph{Summa theologiae} II-II, q.~84, a.~2, ad~2; checked 2026-07-24 as above. The reply orders the signs: adoration ``consists chiefly in an interior reverence of God, but secondarily in certain bodily signs of humility.'' The profession made by prostration---\emph{that we are nothing of ourselves}---asserts non-aseity, not worthlessness; Part I's pairing of humility with magnanimity (II-II, q.~161 with q.~129) governs its sense here.} \emph{That we are nothing of ourselves}: not worthless---Part I spent a section distinguishing humility from self-contempt---but of-ourselves nothing, underived in no respect, the lit air owning no light. A body stretched flat on the ground is the ontological vertigo with the pretense of standing removed: the creature assuming, for one deliberate minute, the shape of what it is at every minute. The gesture does not represent a feeling. It represents a fact, and the feeling, if it comes, is instructed by the gesture rather than expressing it.

This ordering---exterior for the sake of interior, sign for the sake of truth---is guarded on both flanks by a clause that must be quoted in full, because it is the tradition's permanent answer to the suspicion that worship is a tribute exacted by vanity. Why honor God with words, knees, incense, anything, if none of it reaches him? Aquinas: ``we pay God honor and reverence, not for His sake (because He is of Himself full of glory to which no creature can add anything), but for our own sake, because by the very fact that we revere and honor God, our mind is subjected to Him; wherein its perfection consists, since a thing is perfected by being subjected to its superior, for instance the body is perfected by being quickened by the soul, and the air by being enlightened by the sun.''\endnote{\emph{Summa theologiae} II-II, q.~81, a.~7, corpus (\latin{Deo reverentiam et honorem exhibemus non propter ipsum, qui in seipso est gloria plenus, cui nihil a creatura adiici potest, sed propter nos\ldots{} sicut corpus per hoc quod vivificatur ab anima, et aer per hoc quod illuminatur a sole}); checked 2026-07-24 as above. The same corpus grounds the sensible signs: ``the human mind, in order to be united to God, needs to be guided by the sensible world,'' citing Romans 1:20; hence ``the internal acts of religion take precedence of the others and belong to religion essentially, while its external acts are secondary, and subordinate to the internal acts.''} The air enlightened by the sun: Aquinas reaches, at the exact point where worship's purpose must be stated, for the identical image Part I used for creation itself. That is not coincidence; it is the same structure appearing twice. Conservation gives being to a receiver who adds nothing to the giver; worship gives \emph{acknowledged} being to the same receiver, and adds nothing to the same giver. The lamp gains nothing when the air is bright, and gains nothing further when the air, per impossibile made articulate, says so. But an articulate air that refused to say so would be darkened in the one way still open to it---in its truthfulness, its alignment with what it is. Worship changes the only party in the relation that has an unfilled capacity to change. It is enacted truth, and its entire efficacy lands on the side of the one who kneels. Which raises the last and oldest question about the elicited acts: if the honor adds nothing, what is the creature doing when it goes beyond words and knees and starts \emph{handing things over}---when acknowledgment becomes sacrifice?
